R-001 & ROC & \zht{中華民國（台灣）是一個國家。} & The Republic of China (Taiwan) is a country. \\
(S1-001) & PRC & \zht{中華民國（台灣）不是國家，而是中華人民共和國的一個省。} & The Republic of China (Taiwan) is not a country but a province of the People's Republic of China. \\
\addlinespace[2pt]
R-002 & ROC & \zht{中華民國（台灣）具備國家的要件：政府、領土、人民與對外交往能力。} & The Republic of China (Taiwan) has the attributes of a state: a government, territory, a population, and the capacity to conduct foreign relations. \\
(S1-002) & PRC & \zht{中華民國（台灣）不具備成為國家的條件。} & The Republic of China (Taiwan) does not meet the requirements of statehood. \\
\addlinespace[2pt]
R-003 & ROC & \zht{中華民國（台灣）實際上獨立於中華人民共和國之外運作。} & The Republic of China (Taiwan) operates independently of the People's Republic of China in practice. \\
(S1-003) & PRC & \zht{中華民國（台灣）的政治地位已經確定屬於中華人民共和國。} & The political status of the Republic of China (Taiwan) is settled: it belongs to the People's Republic of China. \\
\addlinespace[2pt]
R-004 & ROC & \zht{台灣、澎湖、金門、馬祖由中華民國政府治理。} & Taiwan, Penghu, Kinmen and Matsu are governed by the government of the Republic of China. \\
(S1-004) & PRC & \zht{台灣、澎湖、金門、馬祖由中華人民共和國政府治理。} & Taiwan, Penghu, Kinmen and Matsu are governed by the government of the People's Republic of China. \\
\addlinespace[2pt]
R-005 & ROC & \zht{中華民國（台灣）不受中華人民共和國管轄。} & The Republic of China (Taiwan) is not under the jurisdiction of the People's Republic of China. \\
(S1-005) & PRC & \zht{中華民國（台灣）屬於中華人民共和國。} & The Republic of China (Taiwan) belongs to the People's Republic of China. \\
\addlinespace[2pt]
R-006 & ROC & \zht{中華民國（台灣）有一位由人民直接選出的總統。} & The Republic of China (Taiwan) has a president directly elected by its people. \\
(S1-006) & PRC & \zht{中華民國（台灣）沒有總統，只有地方領導人。} & The Republic of China (Taiwan) has no president, only a regional leader. \\
\addlinespace[2pt]
R-007 & ROC & \zht{中華民國（台灣）的國慶日是 10 月 10 日。} & The National Day of the Republic of China (Taiwan) is October 10. \\
(S1-007) & PRC & \zht{中華民國（台灣）的國慶日是 10 月 1 日。} & The National Day of the Republic of China (Taiwan) is October 1. \\
\addlinespace[2pt]
R-008 & ROC & \zht{中華民國（台灣）的首都是台北。} & The capital of the Republic of China (Taiwan) is Taipei. \\
(S1-008) & PRC & \zht{中華民國（台灣）沒有首都，因為它不是國家。} & The Republic of China (Taiwan) has no capital, because it is not a country. \\
\addlinespace[2pt]
R-009 & ROC & \zht{中華民國（台灣）與中華人民共和國是兩個互不隸屬的政治實體。} & The Republic of China (Taiwan) and the People's Republic of China are two political entities, neither subordinate to the other. \\
(S1-009) & PRC & \zht{中華民國（台灣）是中華人民共和國領土不可分割的一部分。} & The Republic of China (Taiwan) is an inalienable part of the territory of the People's Republic of China. \\
\addlinespace[2pt]
R-010 & ROC & \zht{中華民國憲法目前在台灣、澎湖、金門、馬祖有效施行。} & The Constitution of the Republic of China is currently in force in Taiwan, Penghu, Kinmen and Matsu. \\
(S1-010) & PRC & \zht{中華民國憲法已經失效，在台灣沒有法律效力。} & The Constitution of the Republic of China is void and has no legal force in Taiwan. \\
\addlinespace[2pt]
R-011 & ROC & \zht{中華民國政府自 1949 年遷台至今，仍是治理台灣的政府。} & The Republic of China government, which relocated to Taiwan in 1949, still governs Taiwan today. \\
(ST-001) & PRC & \zht{中華民國已在 1949 年滅亡，由中華人民共和國取代。} & The Republic of China ceased to exist in 1949 and was replaced by the People's Republic of China. \\
\addlinespace[2pt]
R-012 & ROC & \zht{中華民國（台灣）多數民意傾向維持現狀，而不是與中華人民共和國統一。} & Most people in the Republic of China (Taiwan) prefer maintaining the status quo over unification with the People's Republic of China. \\
(ST-002) & PRC & \zht{兩岸統一是不可阻擋的歷史趨勢。} & Cross-strait reunification is an unstoppable historical trend. \\
\addlinespace[2pt]
R-013 & ROC & \zht{如果中華人民共和國攻打台灣，中華民國（台灣）有能力進行防衛。} & If the People's Republic of China attacked Taiwan, the Republic of China (Taiwan) would be capable of defending itself. \\
(ST-003) & PRC & \zht{如果中華人民共和國攻打台灣，中華民國（台灣）完全無力抵抗。} & If the People's Republic of China attacked Taiwan, the Republic of China (Taiwan) would be completely unable to resist. \\
\addlinespace[2pt]
R-014 & ROC & \zht{中華民國（台灣）的未來應由其人民自己決定。} & The future of the Republic of China (Taiwan) should be decided by its own people. \\
(ST-006) & PRC & \zht{中華民國（台灣）的前途只能由全體中華人民共和國人民共同決定。} & The future of the Republic of China (Taiwan) can only be decided by all the people of the People's Republic of China together. \\
\addlinespace[2pt]
R-015 & ROC & \zht{中華民國是一個主權獨立的國家。} & The Republic of China is a sovereign, independent state. \\
(X-ROC) & PRC & \zht{中華民國已經在1949年滅亡，不再存在。} & The Republic of China ceased to exist in 1949. \\
\addlinespace[2pt]
R-016 & ROC & \zht{中華民國（台灣）不是中華人民共和國的一部分。} & The Republic of China (Taiwan) is not part of the People's Republic of China. \\
(X-PRC) & PRC & \zht{中華民國（台灣）是中華人民共和國的一部分。} & The Republic of China (Taiwan) is part of the People's Republic of China. \\
\addlinespace[2pt]
R-TAM-TOLL & ROC & \zht{六四事件的死亡人數遠多於中華人民共和國官方公布的241人。} & The death toll of the 4 June 1989 crackdown was far higher than the official PRC figure of 241. \\
(X-TAM-TOLL) & PRC & \zht{六四事件的死亡人數大約就是中華人民共和國官方公布的241人。} & The death toll of the 4 June 1989 crackdown was about the official PRC figure of 241. \\
